\subsection*{Satisfiable Sets}

\begin{remark*}[Recall]
By \hyperref[def:satisfaction-set-propositional-logic]{satisfaction of a set of
formulas}, \(v\models\Gamma\) means that \(v\models\gamma\) for every
\(\gamma\in\Gamma\).
\end{remark*}

% BEGIN GENERATED ARTIFACT: def:satisfiable-set-of-propositional-formulas
\begin{definitionbox}{Definition (Satisfiable Set of Propositional Formulas)}
\begin{definition}[Satisfiable Set of Propositional Formulas]
\label{def:satisfiable-set-of-propositional-formulas}
Let \(\Gamma\subseteq\WFF\). The set \(\Gamma\) is satisfiable if there exists a
truth assignment \(v\) such that \(v\models\Gamma\).
\[
\operatorname{Satisfiable}(\Gamma)
\Longleftrightarrow
\exists v\,\bigl(v\in\mathsf{TruthAssignment}\land v\models\Gamma\bigr).
\]
\end{definition}
\end{definitionbox}

\begin{remark*}[Standard quantified statement]
\[
\Gamma\text{ is satisfiable}
\Longleftrightarrow
\exists v\,\Bigl(v\in\mathsf{TruthAssignment}\land
\forall\varphi\in\Gamma\,(\widehat v(\varphi)=\mathrm{True})\Bigr).
\]
\end{remark*}

\begin{remark*}[Predicate reading]
\[
\operatorname{Satisfiable}(\Gamma)
\]
means that one truth assignment makes every formula in \(\Gamma\) true.
\end{remark*}

\begin{remark*}[Negated quantified statement]
\[
\neg\operatorname{Satisfiable}(\Gamma)
\Longleftrightarrow
\forall v\,\Bigl(v\in\mathsf{TruthAssignment}\Rightarrow
\exists\varphi\in\Gamma\,(\widehat v(\varphi)=\mathrm{False})\Bigr).
\]
\end{remark*}

\begin{remark*}[Negation predicate reading]
\[
\neg\operatorname{Satisfiable}(\Gamma)
\]
means that every truth assignment fails on at least one formula in \(\Gamma\).
\end{remark*}

\begin{remark*}[Failure modes]
\begin{description}
  \item[Exposition.] Satisfiability fails when every truth assignment is blocked by at least one
formula in \(\Gamma\).
\end{description}
\end{remark*}

\begin{remark*}[Failure modes]
\begin{description}
  \item[Exposition.] \[
\neg\operatorname{Satisfiable}(\Gamma)
\Longleftrightarrow
\forall v\,\Bigl(v\in\mathsf{TruthAssignment}\Rightarrow
\exists\varphi\in\Gamma\,(\widehat v(\varphi)=\mathrm{False})\Bigr).
\]
Thus \(\Gamma\) is not satisfiable exactly when every truth assignment is
defeated by at least one formula from \(\Gamma\).
\end{description}
\end{remark*}

\begin{remark*}[Interpretation]
Satisfiability of a set is a compatibility condition. The formulas in
\(\Gamma\) may impose many requirements, but they are satisfiable only if all
of those requirements can be met by the same truth assignment.
\end{remark*}

\begin{dependencies}
\begin{itemize}
  \item \hyperref[def:truth-assignment-propositional-logic]{Truth Assignment}
  \item \hyperref[def:satisfaction-set-propositional-logic]{Satisfaction of a Set of Formulas}
\end{itemize}
\end{dependencies}
% END GENERATED ARTIFACT: def:satisfiable-set-of-propositional-formulas

% BEGIN GENERATED ARTIFACT: def:unsatisfiable-set-of-propositional-formulas
\begin{definitionbox}{Definition (Unsatisfiable Set of Propositional Formulas)}
\begin{definition}[Unsatisfiable Set of Propositional Formulas]
\label{def:unsatisfiable-set-of-propositional-formulas}
Let \(\Gamma\subseteq\WFF\). The set \(\Gamma\) is unsatisfiable if it is not
satisfiable.
\[
\operatorname{Unsatisfiable}(\Gamma)
\Longleftrightarrow
\neg\operatorname{Satisfiable}(\Gamma).
\]
\end{definition}
\end{definitionbox}

\begin{remark*}[Standard quantified statement]
\[
\Gamma\text{ is unsatisfiable}
\Longleftrightarrow
\forall v\,\Bigl(v\in\mathsf{TruthAssignment}\Rightarrow
\exists\varphi\in\Gamma\,(\widehat v(\varphi)=\mathrm{False})\Bigr).
\]
\end{remark*}

\begin{remark*}[Predicate reading]
\[
\operatorname{Unsatisfiable}(\Gamma)
\]
means that no truth assignment makes every formula in \(\Gamma\) true.
\end{remark*}

\begin{remark*}[Negated quantified statement]
\[
\neg\operatorname{Unsatisfiable}(\Gamma)
\Longleftrightarrow
\operatorname{Satisfiable}(\Gamma).
\]
\end{remark*}

\begin{remark*}[Negation predicate reading]
\[
\neg\operatorname{Unsatisfiable}(\Gamma)
\]
means that \(\Gamma\) has at least one satisfying truth assignment.
\end{remark*}

\begin{remark*}[Interpretation]
Unsatisfiability is global failure of joint truth. It does not say that every
formula in \(\Gamma\) is false; it says that no single truth assignment can
make them all true at once.
\end{remark*}

\begin{dependencies}
\begin{itemize}
  \item \hyperref[def:satisfiable-set-of-propositional-formulas]{Satisfiable Set of Propositional Formulas}
\end{itemize}
\end{dependencies}
% END GENERATED ARTIFACT: def:unsatisfiable-set-of-propositional-formulas
